% Work/MAT -- 1. rocnik -- hodina 40 -- finalni revize 2026-10-30.
% Priklady z uzivatelem dodanych obrazku upraveny a doplneny vlastnimi ulohami.
% Vysledky a podminky kontrolovany nezavisle; zapis vysledku nepredpoklada
% kladnost promennych tam, kde to neni stanoveno.
\documentclass[11pt,a4paper]{article}
% Spolecny styl pro ucitelske pripravy MAT (Work/MAT)
% Verze 1.1, 2026-10-23
\NeedsTeXFormat{LaTeX2e}
\ProvidesPackage{priprava}[2026/10/23 v1.1 Sablona ucitelskych priprav MAT]

% Zakladni typografie a cestina
\RequirePackage[T1]{fontenc}
\RequirePackage[utf8]{inputenc}
\RequirePackage[czech]{babel}
\RequirePackage{lmodern}
\RequirePackage{microtype}
\RequirePackage{amsmath,amssymb}
\RequirePackage{xcolor}
\RequirePackage{enumitem}
\RequirePackage[most]{tcolorbox}
\RequirePackage{tabularx,booktabs,array}
\RequirePackage{fancyhdr}
\RequirePackage{needspace}
\RequirePackage{geometry}
\RequirePackage[hidelinks]{hyperref}

\geometry{a4paper,top=17mm,bottom=18mm,left=23mm,right=23mm}
\definecolor{notegray}{gray}{0.39}
\setlength{\parindent}{0pt}
\setlength{\parskip}{2.5pt plus 1pt minus 0.5pt}
\setlist[itemize]{topsep=2pt,itemsep=1pt,parsep=0pt,leftmargin=1.45em}
\setlist[enumerate]{topsep=2pt,itemsep=2pt,parsep=0pt,leftmargin=1.65em}
\setcounter{secnumdepth}{0}
\renewcommand{\arraystretch}{1.12}

% Zahlavi: vlevo rocnik | globalni tema, vpravo cislo: nazev hodiny
% Zapati zustava zcela prazdne.
\pagestyle{fancy}
\fancyhf{}
\renewcommand{\headrulewidth}{0.35pt}
\renewcommand{\footrulewidth}{0pt}
\setlength{\headheight}{14pt}
\newcommand{\nastavPripravu}[4]{%
  \fancyhead[L]{\small #1\enspace\textbar\enspace #2}%
  \fancyhead[R]{\small #3: #4}%
}

% Minimalisticke cernobile boxy. Kazda definice zvlast.
\tcbset{commonlesson/.style={enhanced,breakable=false,colback=white,colframe=black,
  colbacktitle=white,coltitle=black,boxrule=0.45pt,
  arc=0.8mm,outer arc=0.8mm,left=3mm,right=3mm,top=1.6mm,bottom=1.6mm,
  toptitle=1mm,bottomtitle=0.7mm,fonttitle=\bfseries,
  attach boxed title to top left={xshift=2.5mm,yshift=-1.15mm},
  boxed title style={colback=white,colframe=white,boxrule=0pt,left=1mm,right=1mm}}}
\newtcolorbox{zadaniBox}[1][]{commonlesson,title={Zadání},#1}
\newtcolorbox{prikladBox}[1][]{commonlesson,title={Příklad},#1}
\newtcolorbox{definiceBox}[1][]{commonlesson,title={Definice},#1}

% Spolecny box pro zadani a jeho reseni, bez mezery mezi castmi.
% Pouziti:
%   \begin{prikladReseniBox}
%     Zadani ulohy.
%     \tcblower
%     \begin{align*} ... \end{align*}
%   \end{prikladReseniBox}
% Spodni cast ma jemny sedy podklad a delici linku; lze v ni sazet
% vice radku matematiky. Obe casti maji jen jeden spolecny ramecek.
\tcbset{solutionlesson/.style={
  commonlesson,
  enhanced jigsaw,
  breakable,
  colbacklower=black!3,
  segmentation style={solid,draw=black!25,line width=0.35pt},
  fontlower=\small\color{notegray},
  before lower={\noindent\textbf{Řešení:}\enspace},
  leftlower=3mm,rightlower=3mm
}}
\newtcolorbox{prikladReseniBox}[1][]{solutionlesson,title={Příklad},#1}
\newtcolorbox{zadaniReseniBox}[1][]{solutionlesson,title={Zadání},#1}

% Pomucky pro autora pripravy: sedive poznamky; \reseni pro starsi pripravy
\newcommand{\poznamka}[1]{\par\begingroup\small\color{notegray}\itshape
  \textnormal{\textbf{Poznámka pro učitele.}}\enspace #1\par\endgroup}
\newcommand{\reseni}[1]{\par\begingroup\small\color{notegray}
  \textbf{Očekávané řešení:}\enspace #1\par\endgroup}
\newcommand{\blok}[1]{\needspace{5\baselineskip}\subsection*{#1}}
\newcommand{\vyrok}[1]{\uv{#1}}

\nastavPripravu{1. ročník}{Mocniny a odmocniny}{40}{Test a úvod do výrazů}
\begin{document}

\poznamka{\textbf{45 minut:} 2 min organizace, \textbf{přibližně 20 minut
průběžný test z mocnin a odmocnin} s vlastním zadáním vyučujícího,
3 min odevzdání, 8 min debata \uv{co je výraz},
8 min očekávání nadcházejícího bloku, 4 min závěrečné otázky.
\textbf{Bez dalších výpočtů z mocnin a odmocnin.} Test zde nemá žádné zadání ani řešení.}

\blok{0.--22. minuta: průběžný test}
\begin{zadaniBox}[title={Organizace pro učitele}]
Použij svůj připravený test na \textbf{20 minut čistého času}.
Před zahájením vyjasni povolené pomůcky, bodování, označení variant,
povinnost uvést postup a kontrolu podmínek.
Po dvaceti minutách práci ukonči a vyber.
Zadání, klíč ani konkrétní řešení testu nejsou součástí této přípravy.
\end{zadaniBox}

\blok{22.--25. minuta: uzavření a změna tématu}
Po vybrání písemek řekni, že téma \emph{mocnin a odmocnin} je nyní uzavřeno.
Opravu konkrétních testů prober až po jejich vyhodnocení.
\textbf{Zbytek hodiny není další test ani další procvičování}:
pouze úvodní rozhovor o algebraických výrazech.

\blok{25.--33. minuta: co vůbec nazýváme výrazem?}
\begin{zadaniBox}[title={Diskusní otázky pro třídu}]
\begin{itemize}
\item Co se změní, když místo čísla $3$ napíšeme písmeno $x$?
\item Kde jste už pracovali s proměnnými? (Vzorce pro obsah, obvod, rychlost,
Pythagorova věta, goniometrické poměry.)
\item Je $3x+2$ už hotové číslo, nebo \uv{návod k výpočtu}?
\item Může být jeden a tentýž výraz definovaný pro některá čísla a pro jiná ne?
\end{itemize}
\end{zadaniBox}

\begin{zadaniReseniBox}[title={Příklad 1: výraz, rovnice, nerovnice, nebo nesmyslný zápis?},breakable=false]
Společně na tabuli roztřiď:
\[
\begin{gathered}
3x+2,\quad \frac{x-1}{x+1},\quad \sqrt{x+4},\quad 5(a-b)^2,\\
x^2-4=0,\quad 2x>3,\quad \frac{x+1}{0},\quad 7+\cdot x.
\end{gathered}
\]
\tcblower
První čtyři jsou \textbf{algebraické výrazy} (některé s omezením):
$3x+2$ a $5(a-b)^2$ pro všechna reálná čísla,
$\frac{x-1}{x+1}$ pro $x\ne-1$,
$\sqrt{x+4}$ pro $x\ge-4$.
Zápisy $x^2-4=0$ a $2x>3$ jsou \textbf{rovnice a nerovnice}, nikoli samotné výrazy.
$\frac{x+1}{0}$ není v reálných číslech \emph{nikdy definován};
$7+\cdot x$ není správně zapsaný matematický výraz.
\end{zadaniReseniBox}
\poznamka{Nejde o další formalizaci definic ani zavádění názvosloví
polynomů. Zde jen aktivuj představu, že s výrazy lze pracovat i bez
konkrétně dosazené hodnoty. Rovnice \uv{má rovnítko}, algebraický výraz sám ne.}

\blok{33.--41. minuta: na co se mají žáci těšit?}
\begin{zadaniBox}[title={Přehled budoucího učiva -- motivační diskuse}]
\begin{enumerate}[label=\arabic*)]
\item \textbf{Členy a mnohočleny}: proč $2x^2-3x+1$ není jen změť znaků.
\item \textbf{Sčítání, násobení a dělení výrazů}: jak se z jednoduchých
pravidel stanou rychlé úpravy.
\item \textbf{Vzorce a rozklady}: $(a+b)^2$, rozdíl čtverců, vytýkání;
jak poznat společnou strukturu.
\item \textbf{Lomené výrazy a podmínky}: kdy se smí krátit a co při krácení
nesmíme ztratit.
\item \textbf{Použití ve vzorcích a rovnicích}: zjednodušení výpočtu,
převody vzorců a pozdější řešení rovnic.
\end{enumerate}
\end{zadaniBox}

\poznamka{Doplň jednoduché návrhy z třídy: obvod obdélníku $2(a+b)$,
obsah $ab$, vztah pro práci $W=Fs$. Nyní jde pouze o motivaci,
neprobírej nové definice mnohočlenů ani vzorce pro rozklad.}

\blok{41.--45. minuta: uzavírací otázka do další hodiny}
\begin{zadaniBox}[title={Otevřený závěr bez domácího výpočtu}]
Na tabuli napiš $2x+3x$, $2x\cdot3x$ a $2x+3$.
Zeptej se: \uv{Co z toho byste uměli zjednodušit už teď a kde si
nejsme jistí proč?} Nech studenty jen navrhovat a zdůvodňovat;
\textbf{formální pravidla pro mnohočleny přijdou až příští hodinu}.
\end{zadaniBox}

\poznamka{Původní čtvrtletní písemná práce z 32. hodiny byla
při reorganizaci učiva přesunuta mimo tento blok. Průběžný test
ve 40. hodině ji automaticky nenahrazuje. Termín čtvrtletní písemky
je nutné upravit samostatně v celoročním plánu.}
\end{document}
